\documentclass{article}
\usepackage{icmlko}
\usepackage{graphicx}
\usepackage{booktabs}
\usepackage{amsmath}

\title{조건이 다섯인 줄 알았다 \\
  {\large 하나가 다른 하나의 복사였고, 그 사실이 가장 유망한 짝 스물다섯을 떨어뜨렸다}}

\begin{document}
\maketitle

\begin{abstract}
직전 노트는 개입-결과 짝 105 개를 「식별 5 조건」에 걸어 \textbf{A 등급 0 개}라 판정했고,
그 하나에 논문 제목$\cdot$초록$\cdot$PR$\cdot$원장이 전부 걸려 있었다.
프레시 티처가 규칙을 풀어 보니 \textbf{조건 다섯 중 하나가 다른 하나의 리터럴 복사}였고
둘은 105/105 상수였다. 실제로 등급을 가른 것은 두 칸뿐이며 \textbf{그 둘의 교집합이 구조적으로 0} 이었다.
이 노트는 사전등록을 새로 쓰고 다시 쟀다.
\claimx{「A = 0」은 재현되지 않는다 --- T1 기준 \textbf{A = 21}(T2 를 더하면 23).}
{어느 규칙 수정 하나로도 A 가 0 에 머물면 이 문장은 무너진다.
실측은 C$\to$A 21 이고 \textbf{아래로 간 짝은 0} 이다.}
그리고 옛 판정에서 살아남은 문장 하나를 좁힌다 --- \textbf{막는 것은 W4 하나}이고,
같이 적혀 있던 W8$\cdot$W2 는 애초에 등급 계산에 들어가지 않았다.
\end{abstract}

\section{무엇이 틀렸나}

직전 노트의 등급 규칙(\texttt{runners/ident901.py:257-267})을 풀면 이렇게 된다.

\begin{itemize}
\item \texttt{:255} \textbf{식5 가 식1 의 리터럴 복사}다 --- 105/105 짝에서 두 값이 같다.
      「5 조건」은 실은 \textbf{4 조건}이고 식1 을 두 번 센다
\item \texttt{:228} 식2(약) $=$ 「예」 \textbf{105/105 상수} $\cdot$ \texttt{:251} 식4 $=$ 「예」 \textbf{105/105 상수}
\item 따라서 $A \iff (\text{식1}=\text{예} \wedge \text{식3}=\text{예})$
\end{itemize}

\claimx{그 교집합이 구조적으로 0 이었다.}
{연속형$\cdot$날짜형이 식1 을 통과할 수 있게 되면 교집합이 열린다 --- 실제로 열렸다.}
식3$=$예를 받은 \textbf{25 짝}은 전부 연속형(16) 아니면 날짜형(9)인데,
\texttt{:219}$\cdot$\texttt{:224} 가 그 두 형을 \textbf{정의상} 식1$=$「모른다」로 보낸다.
\textbf{식3(역인과 아님)을 통과한 짝은 전원 식1 을 통과할 수 없었다.}

이것이 왜 나쁜가: 25 짝은 「A 후보가 아니어서」 죽은 게 아니라
\textbf{다섯 조건 중 가장 까다로운 것을 통과한 유일한 짝들}이었다.
규칙 하나가 가장 유망한 것만 골라 떨어뜨렸다.

\section{다시 쟀다 --- 사전등록이 측정보다 앞선다}

\texttt{docs/prereg\_902\_identify.md} 는 커밋 \texttt{b37ccd019}(2026-08-11T04:43:09+09:00)로
들어갔고 측정 산출물 스탬프는 04:51:11$\to$04:51:18 이다.
\texttt{stamp()} 을 \textbf{실행 시작}에서 불러 시작$\neq$끝이 되게 했다(직전 사이클의 자가 여기서 거짓말했다).

사전등록이 미리 못 박은 것 넷:
① 식1 의 \textbf{연속형$\cdot$날짜형 분할}(중앙값)
② \texttt{식5} 폐지 --- \textbf{독립 측정이 불가능하므로 「4 조건이다」라고 적고 조건 수를 줄인다}
③ \texttt{W8}$\cdot$\texttt{W2} 를 \textbf{등급에 안 들어가는 표지}로 명시
④ \textbf{T1 만의 수}를 헤드라인에 쓴다(직전 사전등록이 「T2 는 T1 수에 절대 안 더한다」고 적어 두고 섞었다).

\claimx{옛 등급을 새 코드로 다시 계산해 \textbf{105/105 일치}를 먼저 확인했다.}
{한 짝이라도 어긋나면 대조의 출발점이 다르므로 이동표 전체가 무효다.}

\section{판정}

\begin{figure}[h]\centering
\includegraphics[width=\linewidth]{figs/fig1.pdf}
\caption{(a) T1 기준 등급 분포 --- A 는 $0 \to 21$. (b) 무엇 때문에 바뀌었나.
(c) 무엇이 막았나 --- W4 하나다.}
\end{figure}

\begin{table}[h]\centering\small
\begin{tabular}{lrrr}
\toprule
 & 901 (5 조건) & 902 (4 조건) & 이동 \\
\midrule
A & \textbf{0} & \textbf{21} & C$\to$A 21 \\
B & 17 & 47 & C$\to$B 30 \\
C & 64 & 13 & C$\to$C 13 \\
\bottomrule
\end{tabular}
\end{table}

분모는 \textbf{T1 81}(T2 24 는 더하지 않는다). \textbf{아래로 간 짝은 0} 이다.

바뀐 이유를 분해하면: 변화 없음 30 $\cdot$ 식1 규칙 신설 28 $\cdot$
\textbf{식5 폐지만으로 2} $\cdot$ 둘 다 21.

\claimx{식5 폐지 \emph{하나만}으로 등급이 바뀐 짝이 있다.}
{그 둘의 새 식1 도 「아니오」다 --- 조건이 5 에서 4 로 준 것만으로 움직였다.
즉 식1 을 두 번 세던 것이 그 자리에서 등급을 깎고 있었다.}
팝업 \texttt{capacity\_per\_session}$\cdot$\texttt{unit\_minutes} 둘이다.

\subsection{식3 $=$ 예였던 25 짝의 결말}

\textbf{A 23 $\cdot$ B 2 $\cdot$ C 0.} 티처가 예측한 「최소 9 $\cdot$ 최대 25」가 실측 23 으로 떨어졌다.
연속형 16 중 A 14, 날짜형 9 는 \textbf{전원 A}.

\section{막는 것은 하나로 줄었다}

\begin{figure}[h]\centering
\includegraphics[width=\linewidth]{figs/fig2.pdf}
\caption{(a) 식1 과 식3 의 값 분포 --- 901 에서 두 「예」의 교집합이 0 이었다.
(b) 지평 팔이 찾은 이미 열려 있던 열.}
\end{figure}

\claimx{직전 노트의 「막는 것은 결측이 아니라 W4$\cdot$W8 이다」에서
\textbf{W8 은 빠져야 한다.}}
{W8 이 등급 계산의 어느 조건에라도 들어가면 이 정정은 틀렸다.
실측은 등급이 식1$\cdot$식3 둘로만 매겨진다는 것이다.}
W8 은 짝마다 조건 없이 붙는 표지였고 W2 도 마찬가지다.
남는 것은 \textbf{W4}(처치 시점 없음) 하나이며 이것은 실측이다 ---
식3 $=$ 「모른다」가 T1 55, B 등급이 못 채운 칸의 T1 45 가 이 하나 때문이다.

\claimx{그리고 W8 을 \textbf{실제로 훑었더니} 0 이 아니었다.}
{그 필드들이 배정 사유가 아니라 라벨 신뢰 사유라는 판단은 \textbf{사후 해석}이다 ---
이 문장은 「훑었고 이런 꼴의 필드가 있다」까지만 주장한다.}
「0」은 \textbf{「없다」가 아니라 「훑었고 그 꼴의 필드가 없다」}로 적어야 한다(조항 59).

\section{정직 표지 --- A 23 에 붙인 것}

\begin{itemize}
\item 연속$\cdot$날짜형의 처치/대조는 \textbf{우리가 중앙값으로 만든 것}이지 자연 대조군이 아니다
\item \textbf{날짜형 9 짝은 처치이면서 동시에 결과 창의 원점}이고 판에서 탈추세 변수로 이미 소비된다.
      \claimx{그래도 식3 을 안 바꿨다.}
      {식3 을 조이면 그 9 짝이 다시 떨어진다 --- 그것이 바로 이 노트가 고치려던
      「가장 유망한 것만 골라 떨어뜨리기」의 재발이다.}
\item 팝업 A 14 는 전부 분모 D1 이다 --- 팝업은 축이 \texttt{npz} 라 D3 되짚기가 안 된다
\item 팝업 \texttt{area\_sqm} 은 \textbf{소수 쪽이 정확히 10 $=$ 문턱}이다. 문턱은 직전 사전등록에서
      물려받았고 이 노트가 안 건드렸다
\item 감도: 중앙값 대신 최대균형 분할이면 판정이 갈리는 짝 5 --- 전부 식3 $=$ 「모른다」라
      C$\to$B 로만 움직이고 \textbf{A 수는 안 바뀐다}
\end{itemize}

\section{배선 --- 심은 결함 셋이 전부 발화했다}

W-F 분할 함수 파손 $\cdot$ W-G 둘째 구현이 이웃 키를 읽게 파손
(정상 경로는 105 짝 전량 두 구현 일치) $\cdot$
W-H 영화 키 정규화 되돌리기 $\to$ \texttt{키 406 $\to$ 매칭 0}(직전 노트의 조용한 0 재현).
D3 합 3{,}775 일치.

\section{한계 --- 모른다고 적는 것}

\begin{itemize}
\item \textbf{효과 크기는 여전히 하나도 안 냈다.} A 21 은 \textbf{다음 사이클 후보이지 채택이 아니다}
\item 식3$\cdot$식4 규칙, 문턱 10, 개입 후보 열 목록을 \textbf{안 바꿨다} --- 열이 빠졌는지는 안 봤다
\item 팝업$\cdot$아이돌 D3 되짚기 불가 --- A 14 와 B 다수가 D1 분모 위에 서 있다
\item 사전등록이 이유 라벨을 셋으로 적었는데 \textbf{「둘 다」 라벨을 추가}했다(23 짝).
      감추지 않고 산출물에 적었다
\item 러너를 두 번 돌렸다 --- 두 번째는 \textbf{보고 키만} 추가했고 등급 규칙은 한 글자도 안 바꿨다
      (등급 분포 동일). 산출물에 고지했다
\end{itemize}

\section{곁에서 나온 것}

같은 사이클의 지평 팔이 물었다 --- 「우리가 이미 매일 쌓고 있는 것이 무엇을 열었나」.
답은 둘로 갈렸다.
\claimx{직전 노트가 T1 로 센 열에 대해서는 \textbf{안 열리고, 시간이 해결하지도 않는다.}}
{개입 값이 사는 저장소를 다시 읽는 수집기가 하나라도 생기면 이 문장은 무너진다.
지금은 0 개이고, 그래서 같은 레코드의 두 번째 판이 저장소에 없다.}
\claimx{그런데 \textbf{T1 로 세지 않은 열이 이미 열려 있었다} --- 일별 스크린수$\cdot$상영횟수.}
{그 열이 개입인지는 \textbf{이 노트가 판정하지 않았다.} 배급 결정은 직전 흥행에 반응하므로
동시성이 남는다.}
영화 D3 유보 404 개에서 라벨 창 안에 스크린수 변화가 있는 개체가 398 이다.
직전 노트가 못 본 이유도 나왔다 --- 영화 원천을 \textbf{영화당 한 행짜리 파일 하나}로 잡아서
일별 표가 재고에 한 번도 안 들어갔다.

\end{document}
